\section{Setup and assumptions}

Throughout, \(n,d\geq2\) are integers and \(0<s\leq1\). Unit indices \leanref{sym:i}{\ensuremath{i}} range from \(1\) to \(n\), and coordinate indices \leanref{sym:j}{\ensuremath{j}} and \leanref{sym:ell}{\ensuremath{\ell}} range from \(1\) to \(d\), with \(j<\ell\) for coordinate pairs. A cube point is denoted by \leanref{sym:x}{\ensuremath{x}}. Integrals over cubes use Lebesgue measure, and torus integrals use normalized Haar probability. The notation \(\|\cdot\|_2\) denotes the Euclidean norm for vectors and the \(L^2\) norm for functions, as determined by the argument. Generic positive constants may change between occurrences; a subscript \(s\) records permitted dependence on smoothness. For nonnegative quantities, \(a\lesssim_s b\) means \(a\leq C(s)b\), and \(a\asymp_s b\) means that both comparison inequalities hold.

\subsection{Original-unit sampling}

The full covariate sample \leanref{sym:X}{\ensuremath{X=(X_1,\ldots,X_n)}} records the covariates of all original experimental units. Assignment takes values in the sign space \leanref{sym:Zspace}{\ensuremath{\mathcal Z_n=\{-1,1\}^n}}. An assignment vector \leanref{sym:Z}{\ensuremath{Z=(Z_1,\ldots,Z_n)}} selects treatment through the indicators \leanref{sym:A}{\ensuremath{A_i=(Z_i+1)/2}}. A covariate-dependent assignment law \leanref{sym:pi}{\ensuremath{\pi}} is a Borel probability kernel from \([0,1]^{nd}\) to \(\mathcal P(\mathcal Z_n)\), the probability measures on the finite sign space. Thus the design can use the entire original covariate array when assigning every unit.

The sampling distribution fixes the population measure against which component centering, regularity, and expected imbalance are evaluated.

\begin{assumptionv}{ass:uniform-draw}[Independent uniform covariate sample]
The covariate sample \(X=(X_1,\ldots,X_n)\) satisfies
\[
\operatorname{Law}(X)=\operatorname{Unif}([0,1]^d)^{\otimes n}.
\]
\end{assumptionv}

Independent uniform sampling is the exact uniform special case of the independent sampling and bounded-density condition in \citep[Assumption 2.1]{Cytrynbaum2026Limits}. It supplies a common population distribution for the design criterion and the canonical component decomposition. Independence also separates the randomness of the original sample from the dependence among treatment assignments that the design creates.

\subsection{Canonical components and regularity}

The prognostic function \leanref{sym:m}{\ensuremath{m}} describes the conditional mean of the potential-outcome midpoint. We work with its centered version, whose integral over the cube is zero. Under the product sampling distribution, the canonical main effects \leanref{sym:g1}{\ensuremath{g_j(m)}} and canonical pair effects \leanref{sym:g2}{\ensuremath{g_{j\ell}(m)}} are obtained by marginal integration. The following definition fixes their normalization and the interpretation of almost-everywhere equalities.

\begin{definitionv}{synth_1}[Canonical main and pair effects]
Let \(d\geq2\) and let \(m:[0,1]^d\to\mathbb R\) be a finite-valued Borel function in \(L^2([0,1]^d)\) with Lebesgue integral zero. Its canonical main and pair effects are
\[
g_j(m)(x_j)
=\int_{[0,1]^{d-1}}m(x_j,u_{-j})\,du_{-j},
\qquad 1\leq j\leq d,
\]
and
\[
g_{j\ell}(m)(x_j,x_\ell)
=\int_{[0,1]^{d-2}}
m(x_j,x_\ell,u_{-\{j,\ell\}})\,du_{-\{j,\ell\}}
-g_j(m)(x_j)-g_\ell(m)(x_\ell),
\qquad j<\ell.
\]
The integration arguments fill the remaining coordinates in their original order; a zero-dimensional integral is evaluation. These formulas define equivalence classes almost everywhere. Their centering convention is
\[
\int_0^1g_j(m)(v)\,dv=0,
\qquad
\int_0^1g_{j\ell}(m)(v,x_\ell)\,dv=0,
\qquad
\int_0^1g_{j\ell}(m)(x_j,v)\,dv=0,
\]
where the last two equalities hold for almost every remaining argument. Whenever \(m\) admits an exact order-two decomposition into square-integrable main and pair effects, these centering conditions determine its components uniquely up to Lebesgue null sets.
\end{definitionv}

Coordinatewise centering assigns each pair effect to variation jointly associated with its two coordinates. Together with the zero integral of \(m\), it fixes the intercept and identifies the components in the canonical ANOVA decomposition \citep{Hoeffding1948,Sobol1993}. The specification restricts the prognostic function to the sum of these main and pair effects.

\begin{assumptionv}{ass:order-two}[Exact order-two decomposition]
The prognostic function \(m\) has the decomposition
\[
m(x)=\sum_{j=1}^d g_j(m)(x_j)
+\sum_{j<\ell}g_{j\ell}(m)(x_j,x_\ell)
\]
for Lebesgue-almost every \(x\in[0,1]^d\), where \(g_j(m)\) and \(g_{j\ell}(m)\) are its canonical main and pair effects, respectively.
\end{assumptionv}

The exact order-two canonical ANOVA condition with zero intercept is the specification used here from \citep{Cytrynbaum2026Limits}. It describes a population model in which main effects and pairwise interactions account for the prognostic function exactly. Canonical centering makes the allocation of regularity across components intrinsic to \(m\).

Regularity is measured on the original cube by extension to Euclidean space. For a component dimension \leanref{sym:p}{\ensuremath{p\in\{1,2\}}}, the squared nonperiodic Sobolev restriction norm \leanref{sym:N}{\ensuremath{N_{p,s}(g)^2}} is
\[
N_{p,s}(g)^2
=
\inf_{\substack{G\in L^1(\mathbb R^p)\cap L^2(\mathbb R^p)\\
G|_{[0,1]^p}=g\ {\rm a.e.}}}
\int_{\mathbb R^p}
\left(1+\frac{\|\omega\|_2^2}{p}\right)^s
|\widehat G(\omega)|^2\,d\omega.
\]
Here the Euclidean Fourier transform \leanref{sym:Fourier}{\ensuremath{\widehat G}}, at angular frequency \leanref{sym:omega}{\ensuremath{\omega\in\mathbb R^p}}, uses the convention
\[
\widehat G(\omega)
=(2\pi)^{-p/2}\int_{\mathbb R^p}
e^{-\mathrm i\omega^{\mathsf T}u}G(u)\,du,
\]
and the dimension-normalized weight is exactly \((1+\|\omega\|_2^2/p)^s\). The infimum runs over integrable, square-integrable extensions agreeing almost everywhere on the cube. We impose one common budget across all canonical components.

\begin{assumptionv}{ass:pooled-budget}[Pooled component Sobolev budget]
The canonical main effects \(g_j(m)\) and pair effects \(g_{j\ell}(m)\) satisfy
\[
\sum_{j=1}^d N_{1,s}(g_j(m))^2
+\sum_{j<\ell}N_{2,s}(g_{j\ell}(m))^2
\leq 1,
\]
where \(N_{1,s}\) and \(N_{2,s}\) are the nonperiodic Sobolev restriction norms for one- and two-dimensional components, respectively.
\end{assumptionv}

The pooled component Sobolev budget fixes the dimension-normalized restriction-norm convention used throughout this paper. Pooling bounds the total squared component norm while allowing its distribution across coordinates and pairs to vary. Consequently, the dimension dependence of the criterion reflects the number of available interaction directions under a common regularity budget.

These component conditions define the population class over which the design will be evaluated. Write \leanref{sym:H}{\ensuremath{\mathcal H_{d,s}}} for this centered pooled Sobolev class.

\begin{assumptionv}{ass:fair}[Conditional fairness of assignment]
Under the assignment law \(\pi\), each assignment sign \(Z_i\) has conditional mean zero given the covariate sample \(X\):
\[
\mathbb E_\pi[Z_i\mid X]=0
\qquad\text{for every }i=1,\ldots,n.
\]
\end{assumptionv}

The finite-valued Borel representatives permit evaluation at sampled covariates. Uniform sampling makes changes on Lebesgue null sets immaterial to the expected loss, so the class retains the usual almost-everywhere interpretation of Sobolev functions.

\subsection{Assignment and signing loss}

The assignment family combines covariate-dependent balancing with equal marginal treatment probabilities. Its first requirement concerns each unit's assignment conditional on the complete covariate sample.

\begin{assumptionv}{ass:assignment-independence}[Conditional assignment independence]
A Borel assignment kernel \(\pi\) is outcome-independent if the following holds uniformly over pre-assignment data-generating laws. For every probability law \(\mu_X\) of the full covariate sample and every Markov kernel \(\kappa\) giving the pre-assignment potential-outcome array \(Y\) conditionally on \(X\), the experiment first draws \((X,Y)\) from \(\mu_X\otimes\kappa\) and then draws the assignment signs \(Z\) from \(\pi(X)\). Under this experiment law, \(Z\) and \(Y\) are conditionally independent given \(X\).
\end{assumptionv}

Conditional half-probability marginals are the fairness condition in \citep[Definition 2.2]{Cytrynbaum2026Limits}. Since signs take values in \(\{-1,1\}\), each unit receives treatment with conditional probability \(1/2\). The joint assignment distribution can nevertheless depend on all covariates and can correlate assignments across units to control prognostic imbalance.

The second requirement specifies the information available to the assignment law. Potential outcomes are generated before assignment, and the design draws signs using the covariate array.

\begin{definitionv}{def:sobolev-class}[Pooled Sobolev class \(\mathcal H_{d,s}\)]
The class \(\mathcal H_{d,s}\) consists of functions \(m:[0,1]^d\to\mathbb R\) with the following properties:
\begin{itemize}
\item \textbf{(Measurability.)} \(m\) is finite-valued and Borel measurable.
\item \textbf{(Centering.)} \(m\) belongs to \(L^2([0,1]^d)\) and has Lebesgue integral zero.
\item \textbf{(Component conditions.)} \(m\) satisfies \cref{obj:ass:order-two,obj:ass:pooled-budget} at \((d,s)\).
\end{itemize}
Canonical component equalities and centering are interpreted almost everywhere; arbitrary representatives on null sets give the same risk.
\end{definitionv}

Outcome-independent assignment is the information restriction in \citep[Definition 2.2]{Cytrynbaum2026Limits}. The sequential formulation expresses conditional independence for every pre-assignment outcome law: the design's randomization uses covariates and its own randomness. In conjunction with conditional fairness, it supplies the assignment conditions needed for the causal variance identity below.

Let \leanref{sym:D}{\ensuremath{\mathcal D_{n,d,s}}} denote the resulting admissible design class.

\begin{definitionv}{def:design-class}[Admissible design class $\mathcal D_{n,d,s}$]
The admissible design class is
\[
\mathcal D_{n,d,s}
=
\left\{
\pi:[0,1]^{nd}\to\mathcal P(\mathcal Z_n):
\pi \text{ is a Borel probability kernel satisfying }
\cref{obj:ass:fair,obj:ass:assignment-independence}
\text{ at }(n,d,s)
\right\}.
\]
Here \(\mathcal Z_n\) is the space of length-\(n\) assignment sign vectors, and \(\mathcal P(\mathcal Z_n)\) is its space of probability measures. Each kernel uses the covariates of all original units as its input, with assignment independent of function values, pilot outcomes, and the unknown outcome model.
\end{definitionv}

For a fixed admissible design and a fixed prognostic function, define the scaled signing loss \leanref{sym:loss}{\ensuremath{\mathcal L_{n,d,s}(\pi,m)}} by
\[
\mathcal L_{n,d,s}(\pi,m)
=
\frac4n
\mathbb E_{X,Z}
\left[
\left(\sum_{i=1}^n Z_i m(X_i)\right)^2
\right].
\]
The expectation integrates both the independent uniform sample and the conditional assignment law \(\pi(X)\). The minimax signing risk \leanref{sym:R}{\ensuremath{R_s(n,d)}} optimizes this population loss over admissible designs.

\begin{definitionv}{def:criterion}[Minimax signing risk $R_s(n,d)$]
The minimax scaled signing risk is
\[
R_s(n,d)
=
\inf_{\pi\in\mathcal D_{n,d,s}}
\sup_{m\in\mathcal H_{d,s}}
\frac{4}{n}
\mathbb E_{X,Z}
\left[
\left(\sum_{i=1}^{n}Z_i m(X_i)\right)^2
\right].
\]
Here \(\mathcal D_{n,d,s}\) is the admissible design class, \(\mathcal H_{d,s}\) is the centered prognostic class with pooled order-two Sobolev budget, \(X\) is the independent uniform covariate sample, and \(Z\) is the assignment sign vector drawn under \(\pi(X)\).
\end{definitionv}

The order of optimization is substantive. A design specifies its assignment rule as a function of the covariate array; a fixed population function is then evaluated under that rule, with sampling and assignment averaged inside the loss. The supremum therefore ranges over functions fixed across sample realizations. This order measures how well a single covariate-based design controls population prognostic imbalance uniformly over \(\mathcal H_{d,s}\).

To connect signing loss to treatment-effect estimation, fix the unit-weight Horvitz--Thompson estimator \citep{Horvitz1952}. In the next definition, \(A(a)\) denotes the estimator as a function of a treatment vector \(a\); the subscripted \(A_i\) denotes unit \(i\)'s treatment indicator.

\begin{definitionv}{synth_5}[Unit-weight Horvitz--Thompson estimator]
The unit-weight Horvitz--Thompson estimator \(A\) is defined, for any nonnegative integer \(n\), real potential-outcome schedule \(Y=((Y_i(0),Y_i(1)))_{i=1}^n\), and assignment signs \(Z\in\{-1,1\}^n\), by
\[
A\!\left(\frac{Z+1}{2}\right)
=
\frac{2}{n}\sum_{i=1}^n Z_i
Y_i\!\left(\frac{Z_i+1}{2}\right).
\]
Here \((Z_i+1)/2\) selects treatment when \(Z_i=1\) and control when \(Z_i=-1\). For \(n=0\), the estimator is defined to be zero.
\end{definitionv}

Conditional fairness gives each observed outcome its inverse assignment-probability weight. This estimator remains fixed throughout the design optimization in \cref{obj:def:criterion}.

\subsection{Spectral assignment features}

The assignment construction uses a \leanref{S-4}{reflected and commonly shifted representation} of the original covariates. For an integer \(p\), the period-two torus \leanref{sym:Torus}{\ensuremath{\mathbb T_2^p=[0,2)^p}} has addition modulo two. Its coordinatewise fold \leanref{sym:fold}{\ensuremath{b_p}} is
\[
b_p(y)_h=\min\{y_h,2-y_h\},
\qquad h=1,\ldots,p.
\]
The reflected prognostic function \leanref{sym:F}{\ensuremath{F=m\circ b_d}} is consequently a function on \(\mathbb T_2^d\).

An independent fair bit array \leanref{sym:reflection}{\ensuremath{(b_{ij})}}, with \(b_{ij}\in\{0,1\}\), selects the lifted covariates \leanref{sym:lift}{\ensuremath{\widetilde X_i}} coordinatewise:
\[
\widetilde X_{ij}
=
\begin{cases}
X_{ij},&b_{ij}=0,\\
2-X_{ij},&b_{ij}=1,
\end{cases}
\quad\text{interpreted modulo two}.
\]
Let \leanref{sym:T}{\ensuremath{T}} be a common Haar-uniform shift, independent of the covariates, reflections, and subsequent rounding randomness. The transformed covariate for original unit \(i\) is \leanref{sym:W}{\ensuremath{W_i=(\widetilde X_i+T)\bmod 2}}. Each row used by the assignment procedure still has one entry for each original unit.

For an integer frequency \leanref{sym:k}{\ensuremath{k\in\mathbb Z^d}}, the torus character \leanref{sym:ek}{\ensuremath{e_k(y)=\exp(\mathrm i\pi k^{\mathsf T}y)}} gives the normalized Fourier coefficient \leanref{sym:Fhat}{\ensuremath{\widehat F(k)}} through
\[
\widehat F(k)
=
\int_{\mathbb T_2^d}F(y)\overline{e_k(y)}\,dH(y),
\]
where \(H\) is normalized Haar probability. The frequency complexity \leanref{sym:t}{\ensuremath{t(k)}} ranks frequencies by squared Euclidean magnitude divided by the number of active coordinates. The ordered real features \leanref{sym:features}{\ensuremath{(\phi_h)_{h\geq1}}} convert this ranking into the prescribed rows for assignment.

\begin{definitionv}{synth_3}[Ordered real spectral features and rounding rows]
For \(d\geq2\), write \(\mathbb T_2^d=[0,2)^d\) with addition modulo two. For each \(k\in\mathbb Z^d\) with
\[
1\leq|\operatorname{supp}(k)|\leq2,
\qquad
\operatorname{supp}(k)=\{j:k_j\ne0\},
\]
choose the representative of \(\{k,-k\}\) whose first nonzero coordinate is positive. Order these representatives by increasing
\[
t(k)=\frac{\|k\|_2^2}{|\operatorname{supp}(k)|},
\]
breaking ties lexicographically. For each representative, list the two real features
\[
\sqrt2\cos(\pi k^{\mathsf T}y),
\qquad
\sqrt2\sin(\pi k^{\mathsf T}y),
\qquad y\in\mathbb T_2^d,
\]
consecutively in that order. This defines the ordered sequence \((\phi_h)_{h\geq1}\).

For a reflected and shifted sample \(W_1,\ldots,W_n\), the prescribed real spectral rows are
\[
a_h=\bigl(\phi_h(W_1),\ldots,\phi_h(W_n)\bigr)\in\mathbb R^n,
\qquad h\geq1.
\]
With \(K=\lfloor n/4\rfloor\), the input matrix for ordered boundary rounding is precisely
\[
a=\bigl(\phi_h(W_i)\bigr)_{\substack{1\leq h\leq K\\1\leq i\leq n}}
\in\mathbb R^{K\times n}.
\]
The feature ordering and the choice of rows depend on \(d\) and \(n\) and are independent of the smoothness parameter \(s\).
\end{definitionv}

The support restriction matches the main and pair structure of the prognostic class. Ordering the real features before selecting rows gives a single prescribed balancing matrix for each sample size and dimension. Its independence from \(s\) supplies the common feature input for the smoothness-independent assignment procedure.

\subsection{Outcome laws and causal excess variance}

The causal interpretation uses a single-unit law \(P\) of covariates \(U\) and potential outcomes \(O(0),O(1)\). Define its average treatment effect and conditional prognostic function by
\[
\vartheta(P)=\mathbb E_P[O(1)-O(0)],
\qquad
m_P(u)=
\mathbb E_P\!\left[
\frac{O(1)+O(0)}2\;\middle|\;U=u
\right].
\]
Let \leanref{sym:Q_s}{\ensuremath{\mathcal Q_{d,s}}} consist of these laws with uniform covariates, finite second moments for both potential outcomes, and \(m_P=m\) almost everywhere for some \(m\in\mathcal H_{d,s}\). Thus this outcome-law class fixes the prognostic restriction while allowing square-integrable conditional outcome distributions.

From \(n\) independent copies of \(P\), followed by assignment under \(\pi(U_{1:n})\), the actual estimator is
\[
\widehat T_P
=
\frac2n\sum_{i=1}^n Z_i
O_i\!\left(\frac{Z_i+1}{2}\right).
\]
Under the admissibility conditions, its expectation is \(\vartheta(P)\). The law-specific variance benchmark is
\[
\mathsf V(P)
=
\operatorname{Var}_P\!\left(
\mathbb E_P[O(1)-O(0)\mid U]
\right)
+
2\mathbb E_P\operatorname{Var}_P(O(1)\mid U)
+
2\mathbb E_P\operatorname{Var}_P(O(0)\mid U).
\]
This definition supplies the benchmark directly in terms of the outcome law. It combines variation in the conditional treatment effect with the two conditional outcome variances.

Define the causal excess-variance criterion by
\[
\mathcal V_n(\pi,P)
=
n\operatorname{Var}_{P,\pi}(\widehat T_P)-\mathsf V(P).
\]
The published excess-variance identity gives
\[
\mathcal V_n(\pi,P)
=
\frac4n\mathbb E_{P,\pi}
\left[
\left(\sum_{i=1}^n Z_i m_P(U_i)\right)^2
\right]
=
\mathcal L_{n,d,s}(\pi,m_P)
\]
for \(P\in\mathcal Q_{d,s}\) and \(\pi\in\mathcal D_{n,d,s}\)
\citep[Proposition 2.3, equation (2.4), under Assumption 2.1 and Definition 2.2]{Cytrynbaum2026Limits}.
Hence the signing loss is the scaled actual Horvitz--Thompson variance above a benchmark specific to the outcome law.

A \leanref{S-3}{Gaussian causal completion} supplies a concrete outcome law for every member of the prognostic class. Fix the mean-effect constant \leanref{sym:tau0}{\ensuremath{\tau_0=0.2}} and heterogeneity amplitude \leanref{sym:delta}{\ensuremath{\delta=0.1}}, and define the conditional treatment effect \leanref{sym:tau}{\ensuremath{\tau(x)}} by
\[
\tau(x)=\tau_0+\delta\sqrt2\sin(2\pi x_1).
\]
The Gaussian noise array \leanref{sym:epsilon}{\ensuremath{\epsilon}} has law \(\mathcal N(0,I_{2n})\), independently of \(X\). The pre-assignment potential-outcome schedule \leanref{sym:Y}{\ensuremath{Y}} is generated by
\[
Y_i(0)=m(X_i)-\frac{\tau(X_i)}2+\epsilon_i(0),
\qquad
Y_i(1)=m(X_i)+\frac{\tau(X_i)}2+\epsilon_i(1).
\]
Its conditional mean midpoint is \(m(X_i)\), so its single-unit law belongs to \(\mathcal Q_{d,s}\). The actual estimator \leanref{sym:theta_hat}{\ensuremath{\widehat\theta}}, realized finite-population target \leanref{sym:theta_n}{\ensuremath{\theta_n}}, and sampling-law target \leanref{sym:theta}{\ensuremath{\theta}} are
\[
\widehat\theta
=
\frac2n\sum_{i=1}^n Z_iY_i(A_i),
\qquad
\theta_n
=
\frac1n\sum_{i=1}^n\bigl(Y_i(1)-Y_i(0)\bigr),
\qquad
\theta=\mathbb E[\theta_n]=\tau_0.
\]
For this family, the law-specific benchmark is the constant \leanref{sym:Vstar}{\ensuremath{V^*=\delta^2+4=4.01}}. The published identity therefore expresses \(n\operatorname{Var}(\widehat\theta)-V^*\) as the signing loss of \(m\).

Because the Gaussian completion realizes every \(m\in\mathcal H_{d,s}\), the prognostic image of \(\mathcal Q_{d,s}\) is exactly \(\mathcal H_{d,s}\). It follows from the same identity that
\[
\inf_{\pi\in\mathcal D_{n,d,s}}
\sup_{P\in\mathcal Q_{d,s}}
\mathcal V_n(\pi,P)
=
R_s(n,d).
\]
This equality makes the minimax signing criterion a causal excess-variance criterion over the full uniform-covariate outcome-law class. The following result records the relation between that class and its Gaussian subclass, the corresponding published admissibility conditions, and elementary lower and upper guarantees.

\begin{propositionv}{prop:published-class-inclusion}[Published class inclusion and guarantees]*
Let \(0<s\leq1\), and use the prognostic class, design class, and minimax criterion of \cref{obj:def:sobolev-class,obj:def:design-class,obj:def:criterion}. For each integer \(d\geq2\), let \(\mathcal Q_{d,s}\) be the class of single-unit laws \(P\) of \((U,O(0),O(1))\) such that \(U\) is uniform on \([0,1]^d\), both potential outcomes have finite second moments, and, for some \(m\in\mathcal H_{d,s}\),
\[
m_P(u):=\mathbb E_P\!\left[\frac{O(1)+O(0)}2\;\middle|\;U=u\right]
=m(u)
\quad\text{for almost every }u\in[0,1]^d.
\]
Then the following statements hold.
\begin{itemize}
\item The laws obtained from \(m\in\mathcal H_{d,s}\) by the stated Gaussian causal completion form a proper subset of \(\mathcal Q_{d,s}\), for every integer \(d\geq2\).

\item For every pair of integers \(n,d\geq2\) and every Borel assignment probability kernel \(\pi\) based on the original covariate array, membership in \(\mathcal D_{n,d,s}\) is equivalent to the following published admissibility conditions: under independent uniform covariates \(X\),
\[
\mathbb E_\pi[Z_i\mid X]=0
\quad\text{almost surely, for every }i=1,\ldots,n,
\]
and, for every single-unit law \(P\) whose covariate density is bounded above and below by finite strictly positive constants and whose potential outcomes have finite second moments, the assignment signs are conditionally independent of the full potential-outcome array given the original covariate array in the experiment formed from \(n\) independent copies of \(P\) and the kernel \(\pi\). These are the admissibility conditions of \citet[Assumption 2.1 and Definition 2.2]{Cytrynbaum2026Limits}, specialized to the stated covariate distribution and available information.

\item For every integer \(n\geq2\), the two-dimensional minimax criterion satisfies
\[
R_s(n,2)\geq
\frac{2}{(48+32\pi^2)\,16384^s}\,n^{-s}.
\]

\item For every pair of integers \(n,d\geq2\) and every \(m\in\mathcal H_{d,s}\), independent fair assignment signs, independent of the original covariate sample, give
\[
\frac4n\,\mathbb E\!\left[
\left(\sum_{i=1}^n Z_i m(X_i)\right)^2
\right]
=
4\int_{[0,1]^d}m(x)^2\,dx
=
4\mathbb E[m(X_1)^2]
\leq4,
\]
where \(X_1,\ldots,X_n\) are independent uniform covariate vectors.
\end{itemize}
\end{propositionv}
% CAUSALSMITH-CITED-SCOPE-BEGIN prop:published-class-inclusion
\verificationfootnotetext{\textbf{Formalization scope.} This result uses the cited conclusion \citep{Cytrynbaum2026Limits} (Proposition 2.3 (Excess Variance), equation (2.4), with Assumption 2.1 and Definition 2.2) and \citep{BartlettMendelson2002Complexities} (Theorem 12(4), with Rademacher complexity as in Definition 2). The cited source proof is not formalized here: Lean verifies its use and all remaining steps. If that cited conclusion is false, the portions of this result that depend on it are not certified.}
% CAUSALSMITH-CITED-SCOPE-END prop:published-class-inclusion

\Cref{obj:prop:published-class-inclusion} distinguishes the outcome distributions from the prognostic functions that determine excess variance. The Gaussian family is a proper subclass of the admissible outcome laws, while realizing every prognostic function in the pooled Sobolev class. The admissibility equivalence aligns the original-unit assignment family with the published causal framework. Finally, the two-dimensional lower bound supplies an \(n^{-s}\) benchmark, and independent fair assignment supplies the universal upper bound of \(4\); these comparisons locate the design gains characterized by the main results.
